\section*{Chapter \ref{ch:g12:past-climates} --- Reading Past Climates}

\begin{solution}{exo:g12:past-climates:1}
(a) a \omterm{def:g12:past-climates:proxy}{proxy}; (b) an \omterm{def:g12:past-climates:ice-core}{ice core}; (c) an \omterm{def:g12:past-climates:glacial-cycle}{interglacial}; (d) \omterm{def:g12:past-climates:pollen}{pollen analysis};
(e) the \omterm{def:g12:past-climates:lgm}{Last Glacial Maximum}.
\end{solution}

\begin{solution}{exo:g12:past-climates:2}
(a) Polar ice: the carbon dioxide and methane of the air of the time.
(b) Lake sediments or peat: the vegetation, here trees that need warm
summers. (c) Ocean sediments: the volume of ice on land and the temperature
of the \omterm{def:g10:layered-ocean:thermocline}{deep ocean}. (d) Trees: the growing conditions, year by year.
\end{solution}

\begin{solution}{exo:g12:past-climates:3}
(a) $-50 - (-120) = \qty{70}{m}$ in \num{7000} years: $\qty{70000}{mm}/\num{7000}
= \qty{10}{mm}$ a year. (b) $10/3.7 \approx 2.7$: the mean rise of the
deglaciation was about 2.7 times faster than that of 2006--2018.
\end{solution}

\begin{solution}{exo:g12:past-climates:4}
A century is ten decades: $\qty{2.0}{\degreeCelsius}$ per century. A thousand
years is a hundred decades: $\qty{20}{\degreeCelsius}$ per thousand years.
\end{solution}

\begin{solution}{exo:g12:past-climates:5}
(a) \num{23000} years: the precession of the axis; \num{42000} years: the
tilt of the axis; about \num{100000} years: the shape of the orbit.
(b) $\num{1000000}/\num{41000} \approx 24$ cycles; $\num{1000000}/\num{100000}
= 10$ cycles.
\end{solution}

\begin{solution}{exo:g12:past-climates:6}
(a) $266.6 + 7.83 \times (-9.4) \approx 266.6 - 73.6 = \qty{193}{ppm}$: only
\qty{3}{ppm} above the measured \qty{190}{ppm}. (b) $266.6 + 7.83 \times 5.5
\approx \qty{310}{ppm}$, \qty{11}{ppm} above the highest measured value: the
warmest point lies below the line. (c) $\Delta T = (427 - 266.6)/7.83 \approx
\qty{20}{\degreeCelsius}$. The line only describes how the two quantities
varied together, in natural cycles, within the range of the data
(\qty{-11}{\degreeCelsius} to \qty{+6}{\degreeCelsius}, 174 to
\qty{299}{ppm}); it is not a law, and a value far outside that range cannot
be read from it.
\end{solution}

\begin{solution}{exo:g12:past-climates:7}
(a) Herbs, \qty{84}{\percent}: an open landscape with almost no trees, cold.
(b) Oak (\qty{42}{\percent}) and hazel (\qty{34}{\percent}): a broadleaf
forest, with warm summers. (c) Birch spreads first, then pine; hazel and oak
arrive last. (d) Trees spread only as fast as their seeds travel and their
young trees grow, from the places where they survived the cold; a soil must
also form. The forest can lag behind the climate.
\end{solution}

\begin{solution}{exo:g12:past-climates:8}
(a) \qty{300}{cm}, that is \qty{3000}{mm}, in \num{8000} years: $3000/8000 =
\qty{0.375}{mm}$ a year. (b) $\num{6000} + (470 - 300)/300 \times \num{8000}
\approx \num{6000} + \num{4530} \approx \num{10500}$ years. (c) $10/0.375
\approx 27$ years.
\end{solution}

\begin{solution}{exo:g12:past-climates:9}
(a) $7.16 - 4.11 = 3.05$ million square kilometres, and $3.05/4.11 \approx
\qty{74}{\percent}$ more land. (b) Yes: at the \omterm{def:g12:past-climates:lgm}{Last Glacial Maximum} the sea
stood 125 to \qty{134}{m} below today's, far more than the \qty{36}{m}
that join Borneo to Asia on the map. (c) The map uses today's sea floor, which has not
stayed exactly the same since (sediments have been laid down on it, and the
ground has moved a little); the grid is also coarse, one point every 0.1
degree, and the lowstand is only known to within about ten metres.
\end{solution}

\begin{solution}{exo:g12:past-climates:10}
(a) The value of the \omterm{def:g12:past-climates:lgm}{Last Glacial Maximum}, \qty{4.92}{\permil}, is the
higher. This is expected: the large ice sheets of the time stored water poor
in oxygen-18, which left the ocean, and the shells grown in it, richer in
oxygen-18; the deep water was colder too. (b) So that cold, icy times plot
low and warm times high, as on a temperature curve.
\end{solution}

\begin{solution}{exo:g12:past-climates:11}
(a) A correlation shows that the two records rise and fall together, close
to a straight line; it does not show which caused which, nor that either
did. (b) A warmer climate could have released carbon dioxide into the air;
or a third cause, such as the changes of the Earth's orbit, could have moved
both; or each could have pushed the other. (c) The timing of the changes,
which came first, from a precise dating of the ice and of its air; and an
understanding of how carbon dioxide acts on the Earth's radiation
(\cref{ch:g12:climate-change}).
\end{solution}

\begin{solution}{exo:g12:past-climates:12}
(a) $\num{1300000}/\num{41000} \approx 32$ cycles between 2.6 and 1.3
million years ago; $\num{700000}/\num{98500} \approx 7$ cycles over the last
\num{700000} years. (b) \num{41000} years: the tilt of the axis; about
\num{100000} years: the shape of the orbit. (c) A cycle of \num{41000} years
is sampled about $\num{41000}/\num{2000} \approx 20$ times, enough to draw it;
a change lasting a century falls between two samples, or within one, and
leaves almost no trace.
\end{solution}

\begin{solution}{exo:g12:past-climates:13}
(a) $192/10 \approx 19$ and $198/10 \approx 20$: about 20 times faster.
(b) Against \qty{4}{ppm} per century: $192/4 = 48$ to $198/4 \approx 50$
times; against \qty{42}{ppm}: $192/42 \approx 4.6$ to $198/42 \approx 4.7$
times. Today's rise is about 5 to 50 times faster than the PETM's.
(c) The rise is $2000 - 900 = \qty{1100}{ppm}$: $1100/42 \approx 26$ centuries,
about \num{2600} years at the fastest PETM rate; $1100/195 \approx 5.6$
centuries, under 600 years, at today's rate.
\end{solution}

\begin{solution}{exo:g12:past-climates:14}
(a) Ten times 1 to \qty{1.5}{\degreeCelsius}: 10 to \qty{15}{\degreeCelsius}
per thousand years, that is 0.10 to \qty{0.15}{\degreeCelsius} per decade.
(b) Today's global rate, \qty{0.20}{\degreeCelsius} per decade, is 1.3 to 2
times the fastest warming over Greenland. (c) A single place can change much
more than the average over the globe, where the changes of different regions
partly cancel: a global rate is compared with a global rate.
\end{solution}

\begin{solution}{exo:g12:past-climates:15}
(a) \qty{14}{m} at \qty{12}{m} per thousand years takes $14/12 \approx
1.17$ thousand years, about \num{1170} years: about \num{9800} years ago.
(b) No: the sea had to stay at least \qty{36}{m} below today's to join
Borneo to Asia; once it passed \qty{-36}{m} the land bridge drowned, and
Borneo was an island again. (c) $36/12 = 3$ thousand years, about
\num{3000} years; in fact the rise slowed down after about \num{8200} years
ago, so it took longer.
\end{solution}

\begin{solution}{pb:g12:past-climates:1}
\textbf{1.} The ice itself, through its heavy isotopes, records the
temperature; the air in its bubbles records the carbon dioxide.

\textbf{2.} $5.5 - (-10.6) = \qty{16.1}{\degreeCelsius}$.

\textbf{3.} $299 - 174 = \qty{125}{ppm}$.

\textbf{4.} The points lie close to a straight line: the two records rose
and fell together. It does not tell which caused which.

\textbf{5.} Records of the past are drawn with the oldest times on either
side; read the wrong way, a warming would be taken for a cooling.

\textbf{6.} A \omterm{def:g12:past-climates:glacial-cycle}{glacial period} is a cold interval within an ice age, when the
ice sheets grow far beyond their present size; an \omterm{def:g12:past-climates:glacial-cycle}{interglacial} is a warm
interval between two \omterm{def:g12:past-climates:glacial-cycle}{glacial periods}, with ice sheets close to today's or
smaller.

\textbf{7.} In an \omterm{def:g12:past-climates:glacial-cycle}{interglacial}, which began about \num{11000} years ago.

\textbf{8.} $\num{800000}/\num{98000} \approx 8$: about eight cycles.

\textbf{9.} The shape of the orbit.

\textbf{10.} A rhythm of \num{41000} years, the period of the tilt of the
axis.

\textbf{11.} $-0.4 - (-9.4) = \qty{9.0}{\degreeCelsius}$.

\textbf{12.} $264 - 190 = \qty{74}{ppm}$.

\textbf{13.} \qty{6}{\degreeCelsius}; the Antarctic change, \qty{9.0}{\degreeCelsius},
is larger, about 1.5 times the global one.

\textbf{14.} $\num{18000} - \num{11000} = \num{7000}$ years.

\textbf{15.} $6/7 \approx \qty{0.86}{\degreeCelsius}$ per thousand years.

\textbf{16.} A thousand years is a hundred decades: $0.86/100 \approx
\qty{0.0086}{\degreeCelsius}$ per decade.

\textbf{17.} $0.20 \times 100 = \qty{20}{\degreeCelsius}$ per thousand years.

\textbf{18.} $20/0.86 \approx 23$ times.

\textbf{19.} $20/1.5 \approx 13$ and $20/1 = 20$: 13 to 20 times.

\textbf{20.} Today's global warming is \textbf{about twenty times faster} than
the warming out of the last ice age: 23 times its mean rate, 13 to 20 times
its fastest.
\end{solution}
